\documentclass[a4paper,12pt]{article}
\usepackage[T2A]{fontenc}
\usepackage[utf8]{inputenc}
\usepackage[russian]{babel}
\usepackage{amsmath,amssymb,mathtools}
\usepackage{icomma}
\usepackage{geometry}
\usepackage{microtype}
\usepackage{array,tabularx,booktabs}
\usepackage{hyperref}
\usepackage{tikz}
\usepackage{xcolor}
\usepackage{fancyhdr}
\usepackage{parskip}
\geometry{left=22mm,right=22mm,top=23mm,bottom=23mm}
\definecolor{accent}{HTML}{087F82}
\definecolor{darkaccent}{HTML}{07575A}
\definecolor{textgray}{HTML}{343B3D}
\definecolor{errorred}{HTML}{A32736}
\definecolor{keyblue}{HTML}{1E577D}
\color{textgray}
\hypersetup{colorlinks=true,linkcolor=darkaccent,urlcolor=darkaccent,pdftitle={Ревью домашней работы — 22 ноября 2026},pdfauthor={Лёвин Артём Александрович}}
\setlength{\parindent}{0pt}
\setlength{\parskip}{0.32em}
\renewcommand{\arraystretch}{1.24}
\pagestyle{fancy}
\fancyhf{}
\fancyhead[L]{\footnotesize Отношения и пропорции}
\fancyhead[R]{\footnotesize Ксения Васильченко}
\fancyfoot[C]{\footnotesize\thepage}
\renewcommand{\headrulewidth}{0.35pt}
\renewcommand{\footrulewidth}{0pt}
\setlength{\headheight}{15pt}
\newcommand{\point}[1]{\par\smallskip\noindent{\color{darkaccent}\bfseries #1}\par}
\newcommand{\err}{\textcolor{errorred}{\bfseries Ошибка:}}
\newcommand{\idea}{\textcolor{keyblue}{\bfseries Ключевая идея:}}
\newcommand{\result}{\textcolor{darkaccent}{\bfseries Итог:}}
\newcommand{\tasklabel}[1]{\phantomsection\label{task:#1}}

\begin{document}
\thispagestyle{empty}
\begin{center}
  \vspace*{2.5cm}
  {\color{darkaccent}\large МАТЕМАТИКА\par}
  \vspace{2.1cm}
  {\Huge Ревью домашней работы\par}
  \vspace{1.2cm}
  {\LARGE\color{accent} Отношения, пропорции и проценты\par}
  \vspace{2.0cm}
  {\large Ксения Васильченко\par}
  \vspace{1.0cm}
  {\large 22 ноября 2026 г.\par}
  \vfill
  {\small Лёвин Артём Александрович\par}
  \vspace{0.18cm}
  {\small эксклюзивно для Ксении Васильченко\par}
  \vspace{1.35cm}
  \begin{tikzpicture}
  \draw[accent,line width=0.8pt]
    (0,0) .. controls (1.2,0.55) and (2.9,-0.48) .. (4.2,0)
          .. controls (5.6,0.52) and (7.2,-0.46) .. (8.6,0);
  \end{tikzpicture}
  \vspace{0.35cm}
\end{center}
\newpage

\section*{Сводная таблица}
\textit{В таблице приведены только задания, где требуется исправление или дополнение. Номер задания ведёт к подробному разбору.}

\smallskip
\begingroup\small
\begin{tabularx}{\linewidth}{@{}>{\centering\arraybackslash}p{11mm}p{29mm}X>{\raggedright\arraybackslash}p{28mm}>{\raggedright\arraybackslash}p{26mm}@{}}
\toprule
№ & Тема & Суть замечания & Ваш ответ & Правильный ответ \\
\midrule
\hyperref[task:8]{8} & Пропорция с неизвестным & В последнем преобразовании перепутан порядок вычитания; изменено обозначение неизвестного. & $2{,}5$ (записано с буквой $a$) & $x=2{,}5$ \\
\addlinespace[5pt]
\hyperref[task:10]{10} & Отношение трёх чисел & На предоставленных фотографиях нет решения и ответа. & Не представлен & $12;\ 30;\ 48$, сумма $90$ \\
\bottomrule
\end{tabularx}
\endgroup

\vspace{0.8cm}
{\color{darkaccent}\rule{\linewidth}{0.45pt}}

Оба разбора содержат объяснение причины затруднения, полный школьный способ решения, проверку и одну самостоятельную задачу для закрепления.
\newpage

\thispagestyle{empty}
\null\vfill
\begin{center}
  {\Huge\color{darkaccent} Разбор ошибок}
  \vspace{0.35cm}
  \par\rule{65mm}{0.6pt}
\end{center}
\vfill\null
\newpage

\tasklabel{8}
\section*{Задание 8. Пропорция с неизвестным}
\point{Текст задания}
Найдите $x$ из пропорции
\[
\frac{0{,}8}{x+0{,}5}=\frac{4}{15}.
\]

\point{Ваше решение}
Основная цепочка преобразований записана так:
\begin{align*}
0{,}8:(x+0{,}5)&=4:15,\\
4(x+0{,}5)&=0{,}8\cdot15,\\
4(x+0{,}5)&=12,\\
x+0{,}5&=12:4,\\
x+0{,}5&=3.
\end{align*}
После этого в рукописи вычитание записано в обратном порядке и вместо $x$ используется буква $a$. Последняя строка содержит значение $2{,}5$.

\point{Ваш ответ}
Число $2{,}5$ получено верно, однако обозначение неизвестного в конце изменено.

\par\smallskip
\err{} Неверная запись перехода от $x+0{,}5=3$ к выражению для неизвестного.

\point{Где именно возникает ошибка}
\textbf{Последний правильный шаг:}
\[
x+0{,}5=3.
\]
\textbf{Первый неправильный шаг:} после этой строки в записи меняется местами уменьшаемое и вычитаемое (появляется порядок $0{,}5-3$), а также меняется буква неизвестного. Между тем для нахождения $x$ нужно вычислить именно $3-0{,}5$.

\point{Почему это неверно}
Слева находится сумма неизвестного числа и $0{,}5$. Чтобы оставить неизвестное одно, надо \emph{вычесть $0{,}5$ из обеих частей равенства}:
\[
(x+0{,}5)-0{,}5=3-0{,}5.
\]
Слева два противоположных слагаемых взаимно уничтожаются, остаётся $x$. Справа уменьшаемое равно $3$, а вычитаемое --- $0{,}5$. Порядок вычитания менять нельзя:
\[
3-0{,}5=2{,}5,\qquad 0{,}5-3=-2{,}5.
\]
Это противоположные числа. Если сначала записать неправильную разность, а затем всё равно указать $2{,}5$, получится противоречие между ходом решения и ответом. Также переменная задачи называется $x$, поэтому в конечной записи следует сохранить именно эту букву.

\par\smallskip
\idea{} При переносе слагаемого в другую часть уравнения его знак меняется. То же действие можно понимать как вычитание одного и того же числа из обеих частей уравнения. Важно не менять местами числа в разности.

\point{Исправление от последнего правильного шага}
Продолжим с того места, где преобразования были верными:
\begin{align*}
x+0{,}5&=3,\\
x&=3-0{,}5,\\
x&=\boxed{2{,}5}.
\end{align*}

\point{Полное правильное решение}
В исходной дроби знаменатель не может равняться нулю. Значит,
\[
x+0{,}5\ne0,\qquad x\ne-0{,}5.
\]
Пропорцию можно проверить и решить по основному свойству: произведение крайних членов равно произведению средних. Умножаем обе части равенства на $15(x+0{,}5)$, что допустимо при указанном ограничении:
\begin{align*}
\frac{0{,}8}{x+0{,}5}&=\frac4{15},\\[3pt]
15\cdot0{,}8&=4(x+0{,}5),\\
12&=4(x+0{,}5),\\
3&=x+0{,}5,\\
x&=3-0{,}5,\\
x&=\boxed{2{,}5}.
\end{align*}

\point{Проверка}
Полученное число не нарушает ограничение: $2{,}5\ne-0{,}5$. Подставляем его в исходное равенство:
\[
\frac{0{,}8}{2{,}5+0{,}5}=
\frac{0{,}8}{3}=
\frac{8}{30}=
\frac4{15}.
\]
Левая часть равна правой, значит, найденное значение подходит.

\par\smallskip
\result{} Верный ответ --- $x=2{,}5$. Почти всё решение проведено правильно; исправления требует именно последний переход и обозначение неизвестного.

\point{Задача на закрепление}
Найдите $x$ из пропорции
\[
\frac{1{,}2}{x+0{,}4}=\frac3{10}.
\]

\newpage
\tasklabel{10}
\section*{Задание 10. Три числа в заданном отношении}
\point{Текст задания}
Три числа относятся как $2:5:8$. Самое большое число на $36$ больше самого маленького. Найдите все три числа и их сумму.

\point{Ваше решение}
На предоставленных фотографиях записи решения этого задания нет.

\point{Ваш ответ}
Не представлен.

\par\smallskip
\err{} Задание не завершено: требуется найти не только три числа, но и их сумму.

\point{Где именно решение прерывается}
Последнего правильного вычисления здесь нет, поскольку решение не представлено. Первая необходимая строка, с которой следует начать, --- сравнение числа частей у самого большого и самого маленького числа:
\[
8-2=6\ \text{частей}.
\]

\point{Почему без этого шага задачу не решить}
Отношение $2:5:8$ \emph{не означает}, что сами числа равны $2$, $5$ и $8$. Оно означает, что числа составлены из одинаковых по величине долей: первое содержит две такие доли, второе --- пять, третье --- восемь. Неизвестен размер одной доли.

По условию известна \emph{разность} самого большого и самого маленького числа, а не их сумма. Поэтому число $36$ относится к разности восьми и двух долей, то есть к шести долям. Если сразу разделить $36$ на сумму $2+5+8$, это будет ошибкой: сумма долей здесь не дана условием как реальная сумма чисел.

\par\smallskip
\idea{} Когда несколько чисел заданы отношением, представляйте их как одинаковый множитель, умноженный на каждое число отношения. Если известна разность двух исходных чисел, найдите разность соответствующих частей отношения.

\point{Исправление: как начать решение}
Разность наибольшего и наименьшего коэффициентов:
\[
8-2=6.
\]
Значит, шесть равных долей дают $36$:
\[
36:6=6.
\]
Одна доля равна $6$. Теперь можно вычислить каждое из трёх чисел.

\point{Полное правильное решение. Способ 1: по частям}
\textbf{Шаг 1.} Складывать коэффициенты отношения пока не требуется: сначала используем известную разность. Самое большое число содержит восемь частей, самое маленькое --- две, поэтому их разность содержит
\[
8-2=6\ \text{частей}.
\]
\textbf{Шаг 2.} Эти шесть частей равны $36$. Находим величину одной части:
\[
36:6=6.
\]
\textbf{Шаг 3.} Находим первое число, содержащее две части:
\[
2\cdot6=12.
\]
\textbf{Шаг 4.} Находим второе число, содержащее пять частей:
\[
5\cdot6=30.
\]
\textbf{Шаг 5.} Находим третье число, содержащее восемь частей:
\[
8\cdot6=48.
\]
\textbf{Шаг 6.} Вычисляем сумму всех трёх чисел:
\[
12+30+48=\boxed{90}.
\]

\point{Полное правильное решение. Способ 2: с переменной}
Обозначим размер одной части через $k$. Тогда три числа равны $2k$, $5k$ и $8k$. По условию разность наибольшего и наименьшего составляет $36$:
\begin{align*}
8k-2k&=36,\\
6k&=36,\\
k&=6.
\end{align*}
Следовательно,
\[
2k=12,\qquad5k=30,\qquad8k=48.
\]
Сумма равна $12+30+48=90$.

\point{Проверка}
Проверим оба условия задачи. Числа имеют требуемое отношение:
\[
12:30:48=2:5:8,
\]
потому что каждое из них делится на $6$. Разность наибольшего и наименьшего чисел действительно равна
\[
48-12=36.
\]
Сумма пересчитана отдельно:
\[
12+30+48=90.
\]

\par\smallskip
\result{} Искомые числа --- $12$, $30$, $48$; их сумма --- $90$. Главный шаг --- сопоставить данную разность $36$ с разностью коэффициентов отношения $8-2$.

\point{Задача на закрепление}
Три числа относятся как $3:4:7$. Самое большое число на $32$ больше самого маленького. Найдите все три числа и их сумму.

\vfill
\begin{center}\footnotesize\color{darkaccent}
Лёвин Артём Александрович эксклюзивно для Ксении Васильченко
\end{center}
\end{document}